\documentclass[10pt,UTF8]{article}

\usepackage{geometry}

\usepackage{ctex}

\geometry{a4paper,scale=0.9}
\ctexset{today=old}

\usepackage[bookmarksnumbered=true]{hyperref}  % 生成大纲


\usepackage{makecell}
\usepackage{multirow}

\usepackage{graphicx} %插入图片的宏包
% \usepackage{subfigure}
\usepackage{float} %设置图片浮动位置的宏包

\usepackage{amsmath}
\usepackage{amssymb}
\usepackage{mathtools}

\usepackage{physics}

\usepackage{siunitx}
% 关闭警告
\ExplSyntaxOn
\msg_redirect_name:nnn { siunitx } { physics-pkg } { none }
\ExplSyntaxOff


\usepackage[most]{tcolorbox}


\newenvironment{tipbox}
    {\begin{tcolorbox}[colback=green!5!white,colframe=green!75!black]}
    {\end{tcolorbox}}

\newenvironment{conceptbox}
    {\begin{tcolorbox}[colback=blue!5!white,colframe=blue!75!black]}
    {\end{tcolorbox}}

\newenvironment{thmbox}
    {\begin{tcolorbox}[colback=yellow!5!white,colframe=yellow!75!black]}
    {\end{tcolorbox}}



\newcommand{\mycases}[2][0.8]{
    \left\{\vphantom{\scalebox{#1}{
        $\begin{matrix}#2\end{matrix}$
    }}\right.\!\!\!\!
    \begin{array}{l@{\quad}l}#2\end{array}
}

\newcommand{\perm}[2]{\ensuremath{{}_{#1}\mathrm{P}_{#2}}}  % 排列数
\newcommand{\comb}[2]{\ensuremath{{}_{#1}\mathrm{C}_{#2}}}  % 组合数


\begin{document}

\section{Counting Techniques}

The following enumeration methods are drawn from \emph{multiplication principle}.

\subsection{Summary}

\begin{center}
    \renewcommand{\arraystretch}{2.5}
\begin{tabular}{c|c c|c}
    \hline\hline
    Way to Sample &
    \multicolumn{2}{c|}{Sample You Get} &
    Number of Samples \\
    \hline

    \multirow{2}{*}{
        \makecell{ordered sampling \\ without replacement}        
    } &
    \multicolumn{2}{c|}{a \textbf{permutation of $n$ objects}} &
    $\displaystyle n!$ \\
    \cline{2-4}

    &
    \multicolumn{1}{c|}{\multirow{2}{*}{
        \makecell{an \textbf{ordered sample} \\ \textbf{of size $r$}}
    }} &
    \makecell{a \textbf{permutation of $n$ objects} \\ \textbf{taken $r$ at a time}} &
    $\displaystyle \perm{n}{r}=\frac{n!}{(n-r)!}$ \\
    \cline{1-1} \cline{3-4}

    \makecell{ordered sampling \\ with replacement} &
    \multicolumn{1}{c|}{} &
    &
    $\displaystyle n^r$ \\
    \hline

    \multirow{2}{*}{
        \makecell{unordered sampling \\ without replacement} 
    } &
    \multicolumn{1}{c|}{\multirow{2}{*}{
        \makecell{a \textbf{distinguishable} \\ \textbf{permutation}}
    }} &
    \makecell{
        a \textbf{combination of $n$ objects} \\ \textbf{taken $r$ at a time}
    } &
    $\displaystyle \comb{n}{r}=\binom{n}{r}=\frac{n!}{r!(n-r)!}$ \\
    \cline{3-4}
    
    &
    \multicolumn{1}{c|}{} &
    \makecell{
        a \textbf{distinguishable permutation} \\ \textbf{of objects of $m$ types}
    } &
    $\displaystyle \binom{n}{n_1,n_2,\cdots,n_m}=\frac{n!}{n_1!\cdots n_m!}$ \\
    \hline

    \makecell{unordered sampling \\ with replacement} &
    &
    &
    $\displaystyle \binom{n+r-1}{r}$ \\
    \hline\hline
\end{tabular}
\end{center}


\subsection{Some Notations}

\paragraph{}
The expression $\comb{n}{r}$ can be read as ``$n$ choose $r$''.

\paragraph{}
The numbers $\displaystyle\binom{n}{r}$ are called \textbf{binomial coefficients}.
They arise in the expansion of a binomial.

\paragraph{}
$\displaystyle\binom{n}{n_1,n_2,\cdots,n_s}$ can be called a \textbf{multinomial coefficient}.
\begin{gather*}
    \left(x_1+x_2+\cdots+x_s\right)^n = 
    \sum_{\substack{n_1+n_2+\cdots+n_s=n \\ n_1,n_2,\cdots,n_s\in\mathbb{Z}_{\geqslant 0}}}
    \binom{n}{n_1,n_2,\cdots,n_s}x_1^{n_1}x_2^{n_2}\cdots x_s^{n_s}
\end{gather*}


\subsection{Unordered Sampling with Replacement}

The number of unordered samples of size $r$ selected from $n$ different objects with replacement is equal to
the number of non-negative integer solutions to the equation
\begin{gather*}
    x_1+x_2+\cdots+x_n=r \qquad\text{with } x_i \in \left\{0,1,2,\cdots,r\right\},
    \ i\in\left\{1,2,\cdots,n\right\}
\end{gather*}
where $x_i$ represents the number of times object $i$ is selected.

For each solution $(\hat{x}_1,\hat{x}_2,\cdots,\hat{x}_n)$,
we replace $x_i$ with $\hat{x}_i$ vertical lines in the equation:
\begin{gather*}
    ||| + \ + || + \cdots + || +\  = r
\end{gather*}
Then the number of solutions equals the number of permutations of $r$ vertical lines and $n-1$ plus signs.

Therefore, the number of the original unordered samples is
\begin{gather*}
    \binom{n+r-1}{r} = \binom{n+r-1}{n-1} = \frac{(n+r-1)!}{r!(n-1)!}
\end{gather*}

\end{document}
